\section*{2026 年 10 月 2 日：结论范围更正}
本文保留 9 月 28 日的实验结果作为历史记录。以下说明优先于原正文中较宽泛的解释。
负结果只针对已经测试的损失、复制粘贴与生成配置，不能推出所有数据方法都无效；多种子证据
只适用于实际包含重复训练的实验，并非每项列出的干预。固定标签本身不妨碍施加检测压力，
失败的是当前反事实生成器对真实跨域排序的代理作用。相关性比较只有三个独立架构家族，
不是六个；架构均值的精确置换检验为 $p=0.333$，不能把六点评估解释为强独立架构证据。

代码中的 \texttt{missed\_bg} 指最大 IoU 小于 0.1，不一定完全没有重叠。
mPQ$^+$ 仍与类别有关，$F_c$ 同时受到检测和类型错误影响。发布的验证集最优权重与我们自训
模型的最后检查点，不是完全相同的选点协议。后续审计及 P0--P3 实验（v2，包含合并前发现的
共享边界累加与官方 GT 通道修正）详见仓库 docs 目录的
\texttt{P0\_P3\_RESULTS\_2026-10-02.md}；检测是所研究基线的重要瓶颈，不代表类型预测
对后续改进无关紧要。
\clearpage
